\documentclass[11pt]{article}
\usepackage{mathblatt}
\begin{document}
\weit
\typeout{PROBE-BEGIN flaechen-e1-k2-s4-v2}
\begin{aufgabe}{\detokenize{flaechen-e1-k2-s4-v2}}
\begin{teile}
\teil Probe

\viereck[seiten={$24$ m,$17$ m,$31$ m,$19$ m}]{(0,0)}{(5,0)}{(4.5,3)}{(0.5,2.5)}
\end{teile}
\end{aufgabe}
\typeout{PROBE-END flaechen-e1-k2-s4-v2}
\typeout{PROBE-BEGIN pythagoras-e3-k2-s0-v2}
\begin{aufgabe}{\detokenize{pythagoras-e3-k2-s0-v2}}
\begin{teile}
\teil Probe

\viereck[diagonalen]{(0,0)}{(4.5,0)}{(4.5,2.5)}{(0,2.5)}
\end{teile}
\end{aufgabe}
\typeout{PROBE-END pythagoras-e3-k2-s0-v2}
\typeout{PROBE-BEGIN symmetrie-abbildungen-e2-k2-s3-v2}
\begin{aufgabe}{\detokenize{symmetrie-abbildungen-e2-k2-s3-v2}}
\begin{teile}
\teil Probe

\viereck{(0,0)}{(3,0)}{(3,3)}{(0,3)}
\end{teile}
\end{aufgabe}
\typeout{PROBE-END symmetrie-abbildungen-e2-k2-s3-v2}
\typeout{PROBE-BEGIN terme-e1-k1-s3-v1}
\begin{aufgabe}{\detokenize{terme-e1-k1-s3-v1}}
\begin{teile}
\teil Probe

\viereck[seiten={5,x,5,x}]{(0,0)}{(5,0)}{(5,2.5)}{(0,2.5)}
\end{teile}
\end{aufgabe}
\typeout{PROBE-END terme-e1-k1-s3-v1}
\typeout{PROBE-BEGIN winkel-dreiecke-e2-k2-s9-v2}
\begin{aufgabe}{\detokenize{winkel-dreiecke-e2-k2-s9-v2}}
\begin{teile}
\teil Probe

\viereck[winkel]{(0,0)}{(6,0)}{(4.44,2.5)}{(0.76,2.5)}
\end{teile}
\end{aufgabe}
\typeout{PROBE-END winkel-dreiecke-e2-k2-s9-v2}
\end{document}
